% Reference draft. Every model number comes from generated/numbers.tex through \val{key}, and
% every results table from generated/tab_*.tex; scripts/build_paper.py writes both from
% results/. Rebuild after new runs instead of editing numbers here. Dataset facts written out
% in the text cite their source file in a comment beside them.
\documentclass[11pt,a4paper]{article}

\usepackage[T1]{fontenc}
\usepackage[utf8]{inputenc}
\usepackage{lmodern}
\usepackage[margin=2.4cm]{geometry}
\usepackage{amsmath}
\usepackage{booktabs}
\usepackage{graphicx}
\usepackage{xcolor}
\usepackage{microtype}
\usepackage[numbers,sort&compress]{natbib}
\usepackage[colorlinks=true,linkcolor=blue!50!black,citecolor=blue!50!black,urlcolor=blue!50!black]{hyperref}

% Generated numbers: \val{key} prints the value, or the key in red when no result exists yet.
\newcommand{\nwrdset}[2]{\expandafter\def\csname nwrd@#1\endcsname{#2}}
\newcommand{\val}[1]{\ifcsname nwrd@#1\endcsname\csname nwrd@#1\endcsname\else\textcolor{red}{\textbf{[\detokenize{#1}]}}\fi}
\newcommand{\pending}[1]{\textcolor{red}{\textbf{[Pending: #1]}}}
\newcommand{\pendingtable}[1]{\fbox{\parbox{0.8\linewidth}{\centering\vspace{4pt}\textcolor{red}{\textbf{Pending: #1}}\vspace{4pt}}}}
\newcommand{\note}[1]{\textcolor{blue!60!black}{\textit{[Note to authors: #1]}}}
\newcommand{\figureorpending}[2]{\IfFileExists{figures/#1}{\includegraphics[width=#2]{figures/#1}}{\pendingtable{figure #1}}}
% Written by scripts/build_paper.py from results/: rebuild instead of editing.
\nwrdset{128.canet-unet.ci}{[$-0.0154$, $+0.0177$]}
\nwrdset{128.canet-unet.diou}{$+0.0062$}
\nwrdset{128.canet-unet.p}{0.492}
\nwrdset{128.canet-unet.ttp}{0.847}
\nwrdset{128.canet-unet.wins}{2 of 3}
\nwrdset{128.canet.dice.pm}{0.5788\,$\pm$\,0.0129}
\nwrdset{128.canet.iou}{0.4073}
\nwrdset{128.canet.iou.pm}{0.4073\,$\pm$\,0.0127}
\nwrdset{128.canet.iouhalf.pm}{0.3673\,$\pm$\,0.0111}
\nwrdset{128.canet.pp.f1}{0.517}
\nwrdset{128.canet.pp.f1.vst3}{$-0.040$}
\nwrdset{128.canet.pp.f1rust}{0.638}
\nwrdset{128.canet.pp.f1rust.vst4}{$+0.074$}
\nwrdset{128.canet.pp.f1thr}{0.580}
\nwrdset{128.canet.pp.iourust}{0.468}
\nwrdset{128.canet.pp.iourust.vst4}{$+0.030$}
\nwrdset{128.canet.pp.iourustthr}{0.475}
\nwrdset{128.canet.pp.prec}{0.408}
\nwrdset{128.canet.pp.precrust}{0.559}
\nwrdset{128.canet.pp.rec}{0.735}
\nwrdset{128.canet.pp.recrust}{0.743}
\nwrdset{128.canet.prec.pm}{0.5448\,$\pm$\,0.0137}
\nwrdset{128.canet.rec.pm}{0.6173\,$\pm$\,0.0119}
\nwrdset{128.canet.seeds}{3}
\nwrdset{128.canet.thr}{0.91}
\nwrdset{128.canet.val.iou.pm}{0.4006\,$\pm$\,0.0088}
\nwrdset{128.dlr-canet.ci}{[$-0.0236$, $+0.0298$]}
\nwrdset{128.dlr-canet.diou}{$+0.0024$}
\nwrdset{128.dlr-canet.p}{0.828}
\nwrdset{128.dlr-canet.ttp}{0.437}
\nwrdset{128.dlr-canet.wins}{2 of 3}
\nwrdset{128.dlr-unet.ci}{[$-0.0225$, $+0.0322$]}
\nwrdset{128.dlr-unet.diou}{$+0.0086$}
\nwrdset{128.dlr-unet.p}{0.528}
\nwrdset{128.dlr-unet.ttp}{0.793}
\nwrdset{128.dlr-unet.wins}{2 of 3}
\nwrdset{128.dlr.dice.pm}{0.5812\,$\pm$\,0.0151}
\nwrdset{128.dlr.iou}{0.4098}
\nwrdset{128.dlr.iou.pm}{0.4098\,$\pm$\,0.0150}
\nwrdset{128.dlr.iouhalf.pm}{0.3784\,$\pm$\,0.0098}
\nwrdset{128.dlr.pp.f1}{0.535}
\nwrdset{128.dlr.pp.f1.vst3}{$-0.022$}
\nwrdset{128.dlr.pp.f1rust}{0.653}
\nwrdset{128.dlr.pp.f1rust.vst4}{$+0.089$}
\nwrdset{128.dlr.pp.f1thr}{0.591}
\nwrdset{128.dlr.pp.iourust}{0.485}
\nwrdset{128.dlr.pp.iourust.vst4}{$+0.047$}
\nwrdset{128.dlr.pp.iourustthr}{0.490}
\nwrdset{128.dlr.pp.prec}{0.433}
\nwrdset{128.dlr.pp.precrust}{0.578}
\nwrdset{128.dlr.pp.rec}{0.735}
\nwrdset{128.dlr.pp.recrust}{0.752}
\nwrdset{128.dlr.prec.pm}{0.5255\,$\pm$\,0.0196}
\nwrdset{128.dlr.rec.pm}{0.6504\,$\pm$\,0.0101}
\nwrdset{128.dlr.seeds}{3}
\nwrdset{128.dlr.thr}{0.89}
\nwrdset{128.dlr.val.iou.pm}{0.4163\,$\pm$\,0.0075}
\nwrdset{128.unet.dice.pm}{0.5720\,$\pm$\,0.0367}
\nwrdset{128.unet.iou}{0.4011}
\nwrdset{128.unet.iou.pm}{0.4011\,$\pm$\,0.0365}
\nwrdset{128.unet.iouhalf.pm}{0.3426\,$\pm$\,0.0403}
\nwrdset{128.unet.pp.f1}{0.502}
\nwrdset{128.unet.pp.f1.vst3}{$-0.055$}
\nwrdset{128.unet.pp.f1rust}{0.621}
\nwrdset{128.unet.pp.f1rust.vst4}{$+0.057$}
\nwrdset{128.unet.pp.f1thr}{0.577}
\nwrdset{128.unet.pp.iourust}{0.451}
\nwrdset{128.unet.pp.iourust.vst4}{$+0.013$}
\nwrdset{128.unet.pp.iourustthr}{0.481}
\nwrdset{128.unet.pp.prec}{0.393}
\nwrdset{128.unet.pp.precrust}{0.518}
\nwrdset{128.unet.pp.rec}{0.769}
\nwrdset{128.unet.pp.recrust}{0.782}
\nwrdset{128.unet.prec.pm}{0.4914\,$\pm$\,0.0595}
\nwrdset{128.unet.rec.pm}{0.6888\,$\pm$\,0.0103}
\nwrdset{128.unet.seeds}{3}
\nwrdset{128.unet.thr}{0.88}
\nwrdset{128.unet.val.iou.pm}{0.3656\,$\pm$\,0.0209}
\nwrdset{512.canet-unet.ci}{[$-0.0006$, $+0.0612$]}
\nwrdset{512.canet-unet.diou}{$+0.0252$}
\nwrdset{512.canet-unet.p}{0.062}
\nwrdset{512.canet-unet.ttp}{0.138}
\nwrdset{512.canet-unet.wins}{3 of 3}
\nwrdset{512.canet.dice.pm}{0.5950\,$\pm$\,0.0070}
\nwrdset{512.canet.iou}{0.4235}
\nwrdset{512.canet.iou.pm}{0.4235\,$\pm$\,0.0072}
\nwrdset{512.canet.iouhalf.pm}{0.4027\,$\pm$\,0.0113}
\nwrdset{512.canet.pp.f1}{0.556}
\nwrdset{512.canet.pp.f1.vst3}{$-0.001$}
\nwrdset{512.canet.pp.f1rust}{0.648}
\nwrdset{512.canet.pp.f1rust.vst4}{$+0.084$}
\nwrdset{512.canet.pp.f1thr}{0.603}
\nwrdset{512.canet.pp.iourust}{0.479}
\nwrdset{512.canet.pp.iourust.vst4}{$+0.041$}
\nwrdset{512.canet.pp.iourustthr}{0.476}
\nwrdset{512.canet.pp.prec}{0.476}
\nwrdset{512.canet.pp.precrust}{0.592}
\nwrdset{512.canet.pp.rec}{0.678}
\nwrdset{512.canet.pp.recrust}{0.716}
\nwrdset{512.canet.prec.pm}{0.5779\,$\pm$\,0.0245}
\nwrdset{512.canet.rec.pm}{0.6140\,$\pm$\,0.0120}
\nwrdset{512.canet.seeds}{3}
\nwrdset{512.canet.thr}{0.92}
\nwrdset{512.canet.val.iou.pm}{0.4359\,$\pm$\,0.0061}
\nwrdset{512.dlr-canet.ci}{[$+0.0127$, $+0.0476$]}
\nwrdset{512.dlr-canet.diou}{$+0.0337$}
\nwrdset{512.dlr-canet.p}{0.002}
\nwrdset{512.dlr-canet.ttp}{0.010}
\nwrdset{512.dlr-canet.wins}{3 of 3}
\nwrdset{512.dlr-unet.ci}{[$+0.0415$, $+0.0797$]}
\nwrdset{512.dlr-unet.diou}{$+0.0589$}
\nwrdset{512.dlr-unet.p}{$<$0.001}
\nwrdset{512.dlr-unet.ttp}{0.037}
\nwrdset{512.dlr-unet.wins}{3 of 3}
\nwrdset{512.dlr.dice.pm}{0.6275\,$\pm$\,0.0079}
\nwrdset{512.dlr.iou}{0.4572}
\nwrdset{512.dlr.iou.pm}{0.4572\,$\pm$\,0.0084}
\nwrdset{512.dlr.iouhalf.pm}{0.4423\,$\pm$\,0.0015}
\nwrdset{512.dlr.pp.f1}{0.592}
\nwrdset{512.dlr.pp.f1.vst3}{$+0.035$}
\nwrdset{512.dlr.pp.f1rust}{0.677}
\nwrdset{512.dlr.pp.f1rust.vst4}{$+0.113$}
\nwrdset{512.dlr.pp.f1thr}{0.636}
\nwrdset{512.dlr.pp.iourust}{0.512}
\nwrdset{512.dlr.pp.iourust.vst4}{$+0.074$}
\nwrdset{512.dlr.pp.iourustthr}{0.507}
\nwrdset{512.dlr.pp.prec}{0.519}
\nwrdset{512.dlr.pp.precrust}{0.631}
\nwrdset{512.dlr.pp.rec}{0.700}
\nwrdset{512.dlr.pp.recrust}{0.732}
\nwrdset{512.dlr.prec.pm}{0.6078\,$\pm$\,0.0419}
\nwrdset{512.dlr.rec.pm}{0.6516\,$\pm$\,0.0319}
\nwrdset{512.dlr.seeds}{3}
\nwrdset{512.dlr.thr}{0.82}
\nwrdset{512.dlr.val.iou.pm}{0.4452\,$\pm$\,0.0090}
\nwrdset{512.unet.dice.pm}{0.5697\,$\pm$\,0.0121}
\nwrdset{512.unet.iou}{0.3983}
\nwrdset{512.unet.iou.pm}{0.3983\,$\pm$\,0.0119}
\nwrdset{512.unet.iouhalf.pm}{0.2900\,$\pm$\,0.0176}
\nwrdset{512.unet.pp.f1}{0.442}
\nwrdset{512.unet.pp.f1.vst3}{$-0.115$}
\nwrdset{512.unet.pp.f1rust}{0.605}
\nwrdset{512.unet.pp.f1rust.vst4}{$+0.041$}
\nwrdset{512.unet.pp.f1thr}{0.579}
\nwrdset{512.unet.pp.iourust}{0.434}
\nwrdset{512.unet.pp.iourust.vst4}{$-0.004$}
\nwrdset{512.unet.pp.iourustthr}{0.480}
\nwrdset{512.unet.pp.prec}{0.319}
\nwrdset{512.unet.pp.precrust}{0.475}
\nwrdset{512.unet.pp.rec}{0.827}
\nwrdset{512.unet.pp.recrust}{0.834}
\nwrdset{512.unet.prec.pm}{0.5087\,$\pm$\,0.0469}
\nwrdset{512.unet.rec.pm}{0.6543\,$\pm$\,0.0535}
\nwrdset{512.unet.seeds}{3}
\nwrdset{512.unet.thr}{0.96}
\nwrdset{512.unet.val.iou.pm}{0.3454\,$\pm$\,0.0136}
\nwrdset{data.twins.max}{0.63}
\nwrdset{data.twins.median}{0.50}
\nwrdset{data.twins.min}{0.40}
\nwrdset{data.twins.n}{10}
\nwrdset{paper.t3.f1}{0.557}
\nwrdset{paper.t3.precision}{0.506}
\nwrdset{paper.t3.recall}{0.624}
\nwrdset{paper.t4.f1}{0.564}
\nwrdset{paper.t4.iou}{0.438}
\nwrdset{paper.t4.precision}{0.593}
\nwrdset{paper.t4.recall}{0.552}


\title{A pretrained DeepLabV3+ outperforms the published UNet with adaptive patching on the
NWRD wheat stripe rust dataset: a leakage-free re-benchmark and architectural ablation}
\author{[Authors]\\[2pt]\small [Affiliations]}
\date{Reference draft, \today}

\begin{document}
\maketitle

\begin{center}
\fbox{\parbox{0.92\linewidth}{\small\textbf{How to read this draft.} Numbers in
\textcolor{red}{\textbf{red}} are waiting for runs (the 128px scale and the ablation); blue notes
are for the authors and must be removed before submission. Every other number is generated from
the study's result files and is regenerated, not retyped, when results change.}}
\end{center}

\begin{abstract}
Pixel-level segmentation of wheat stripe rust in field photographs is difficult: rust covers well
under one percent of the pixels, and the only public dataset with pixel annotations, NWRD, holds
100 photographs. Its authors proposed a UNet trained with adaptive patching with feedback (APF)
and reported an F1 score of \val{paper.t3.f1} on their test patches. We re-examine this benchmark.
An audit of the public release shows that it contains 90 distinct scenes rather than 100, that
duplicated photographs sit on both sides of the published train/test and train/validation
boundaries, that the two independent masks of a duplicated photograph agree only at an IoU of
\val{data.twins.min}--\val{data.twins.max}, and that one validation mask is misregistered by
8~pixels. Under a leakage-free protocol that scores every pixel of the 10 test photographs exactly
once, we train the NWRD paper's UNet, an ImageNet-pretrained DeepLabV3+ with a ResNet-50 encoder,
and CANet-B4, a context-aggregation network, with three seeds each. At 512-pixel windows,
DeepLabV3+ reaches a test IoU of \val{512.dlr.iou.pm}, against \val{512.unet.iou.pm} for the UNet
($\Delta$IoU \val{512.dlr-unet.diou}, 95\% CI \val{512.dlr-unet.ci} over test photographs) and
\val{512.canet.iou.pm} for CANet-B4. On the paper's own 2,737 test patches it reaches F1
\val{512.dlr.pp.f1} against the published \val{paper.t3.f1}, although, unlike the published model,
it never trained on a duplicate of a test photograph. \pending{one sentence on the 128px scale and
one on which architectural change the ablation identifies.}
\end{abstract}

\noindent\textbf{Keywords:} wheat stripe rust; semantic segmentation; DeepLabV3+; benchmark;
data leakage; precision agriculture

% =================================================================================================
\section{Introduction}
\label{sec:intro}

Stripe (yellow) rust, caused by \emph{Puccinia striiformis} f.\,sp.\ \emph{tritici}, is among the
most damaging diseases of wheat. Measuring how much of a leaf surface is infected, rather than only
whether a plant is infected, requires pixel-level segmentation, and pixel-level annotation is
expensive. The NWRD dataset of Anwar et al.~\cite{anwar2023nwrd,anwar2023zenodo} is, to our
knowledge, the first public multi-leaf, field-condition segmentation dataset for wheat rust: 100
high-resolution photographs with binary masks. With it, the authors proposed a training pipeline
in which a plain UNet~\cite{ronneberger2015unet} is trained on patches selected by \emph{adaptive
patching with feedback} (APF), which adds patches around the model's own false positives every few
epochs, with a focal loss~\cite{lin2017focal}. They reported precision \val{paper.t3.precision},
recall \val{paper.t3.recall} and F1 \val{paper.t3.f1} on all test patches, and F1 \val{paper.t4.f1} and IoU \val{paper.t4.iou} in a
second setting that also outperformed an Octave-UNet~\cite{anwar2023nwrd}.

Two questions follow for anyone building on this benchmark. First, does a standard, widely
available segmentation architecture with an ImageNet-pretrained encoder do better than the
published model on the same test photographs? Second, if it does, which architectural difference
is responsible? Going from the paper's UNet to DeepLabV3+~\cite{chen2018deeplabv3plus} with a
ResNet-50 encoder~\cite{he2016resnet} changes three things at once: the encoder, its pretraining on
ImageNet~\cite{deng2009imagenet}, and the context module and decoder (atrous spatial pyramid pooling
in place of the UNet's symmetric skip decoder).

Answering either question requires a sound evaluation. We found that the public release contains
repeated photographs across the published splits, so that the published test scores include a
leak, and that the two masks of each repeated photograph disagree substantially, which bounds any
model's attainable IoU.

\paragraph{Contributions.}
\begin{enumerate}
  \item An audit of the public NWRD release: repeated photographs and their cross-split leakage,
  the agreement between independent annotations of the same photograph, and a misregistered mask
  (Section~\ref{sec:data}).
  \item A leakage-free evaluation protocol on the paper's own test photographs that scores every
  pixel exactly once, chooses decision thresholds on validation only, and quantifies uncertainty by
  resampling test photographs (Section~\ref{sec:methods}).
  \item A three-seed comparison of the paper's UNet, DeepLabV3+ R50 and CANet-B4 at two scales, and
  against the published model on its own test patches (Section~\ref{sec:results}).
  \item An ablation that changes one factor at a time between the UNet and DeepLabV3+ (pretraining,
  context module, encoder), identifying which change accounts for the gain
  (Section~\ref{sec:ablation}).
\end{enumerate}

\note{The study was pre-registered (2026-10-22) with CANet-B4 as the candidate model and UNet and
DeepLabV3+ R50 as baselines; the DeepLabV3+-centred framing was chosen after the 512px results. Say
so in the paper (Section~\ref{sec:stats} does), and report CANet-B4's pre-registered verdicts: a
reviewer who learns of the pre-registration and does not find them will read the change of framing
as selective reporting. DeepLabV3+ R50 is the standard implementation of
segmentation-models-pytorch~\cite{iakubovskii2019smp}, so the contribution is the evaluation and
the ablation, not a new architecture.}

% =================================================================================================
\section{Related work}
\label{sec:related}

\paragraph{Wheat rust datasets and the NWRD benchmark.} \note{Add the classification and detection
datasets for wheat rust the target venue expects (image-level datasets, field and leaf-level
collections), and position NWRD as the pixel-level one.} Anwar et al.~\cite{anwar2023nwrd} compared
grid patching, adaptive patching and APF at strides 32 and 128, reporting that APF reaches a
comparable F1 with far shorter training, and compared their UNet with Octave-UNet (F1 0.564 against
0.529, IoU 0.438 against 0.316, one run each).

\paragraph{Segmentation architectures.} UNet~\cite{ronneberger2015unet} pairs a contracting encoder
with a symmetric decoder joined by skip connections at every resolution. DeepLabV3+
\cite{chen2018deeplabv3plus} applies atrous spatial pyramid pooling (ASPP) to the deepest features
of a pretrained backbone and fuses the result with one set of low-level features in a light decoder.
EfficientNet~\cite{tan2019efficientnet} scales depth, width and resolution jointly and is a common
alternative backbone.

\paragraph{Class imbalance.} Rust pixels are rare. Common remedies are the focal
loss~\cite{lin2017focal}, which the NWRD paper used, the Dice loss~\cite{milletari2016vnet} and a
positively weighted binary cross-entropy; we use a fixed combination of the last two for every
model (Section~\ref{sec:training}).

% =================================================================================================
\section{The public NWRD release}
\label{sec:data}

% Source for the numbers in this section: docs/dataset.md and results/dataset/ (summary.md,
% photographs.json, duplicates.json), written by scripts/describe_data.py.

\subsection{Contents}
The release (Zenodo, doi:10.5281/zenodo.21672928, version 1.0, CC~BY~4.0~\cite{anwar2023zenodo})
holds 100 JPEG photographs with 8-bit binary masks, in a training folder (90) and a test folder
(10). The photographs come in 10 sizes, mostly $4608\times3456$ (Huawei PAR-LX1M, either
orientation) and $6016\times4000$ (Nikon D3200); seven are cropped to unique sizes. Every
photograph contains some rust, but rust covers only 0.57\% of all pixels (0.63\% of the test
pixels). We verified every file against its MD5 checksum before every run.

\subsection{Repeated photographs and annotation agreement}
The release holds 90 distinct scenes, not 100. Nine pairs of files decode to identical pixels (same
camera and capture second), and photograph 149 is photograph 99 cropped and re-encoded. Each copy
carries its own mask, and the two masks of a pair agree only at IoU \val{data.twins.min} to
\val{data.twins.max} (median \val{data.twins.median}, \val{data.twins.n} pairs;
Table~\ref{tab:twins}): the same rust, annotated twice, differently. Under the paper's split, two
test photographs (80 and 89, 16.9\% of the test pixels and 9.7\% of the test rust) and two
validation photographs (87 and 149) have their twin in the training set. The published test scores
therefore include a leak. We remove the four training twins (79, 86, 96 and 99) and keep both copies
of the six pairs that lie entirely inside the training set.

\begin{table}[t]
\centering
\caption{Repeated photographs in the public NWRD release and the IoU between their two masks.
Pairs that cross the NWRD paper's splits come first.}
\label{tab:twins}
\small
\begin{tabular}{lcc}
\toprule
Pair & The paper's splits & Mask IoU \\
\midrule
79 / 80 & train / test & 0.509 \\
86 / 89 & train / test & 0.426 \\
87 / 96 & val / train & 0.475 \\
99 / 149 & train / val & 0.583 \\
15 / 32 & train / train & 0.598 \\
69 / 141 & train / train & 0.496 \\
78 / 81 & train / train & 0.522 \\
83 / 128 & train / train & 0.496 \\
91 / 98 & train / train & 0.402 \\
100 / 137 & train / train & 0.632 \\
\bottomrule
\end{tabular}

\end{table}

The agreement between two annotations of one photograph is a useful reference for model scores: a
model's IoU against one annotation cannot be expected to exceed, by much, the IoU between two
annotations of the same scene.

\subsection{A misregistered validation mask}
The mask of photograph 101, one of the paper's nine validation photographs, is 19 pixels narrower
and 21 pixels shorter than the photograph. Searching placements up to 24 pixels past the
photograph's edges, the mask lines up with the rust when shifted 8 pixels left and 1 pixel up from
the photograph's corner, partly outside the photograph, and not at the corner where the authors'
loader placed it (Figure~\ref{fig:mask101}). The six other cropped photographs fit best within
3~pixels of their own corner. Because the mask's right and bottom edges are unknown, we leave
photograph 101 out of every split.

\begin{figure}[t]
\centering
\figureorpending{mask_101.png}{0.85\linewidth}
\caption{Photograph 101 around its three largest annotations (rows). Left: the photograph.
Middle: the mask's outline placed at the photograph's top-left corner (cyan), as the authors' loader
placed it, about 8 pixels right of the rust. Right: shifted 8 pixels left and 1 pixel up (magenta),
where it follows the rust.}
\label{fig:mask101}
\end{figure}

\subsection{The paper's split and patches, reproduced}
From the authors' code\footnote{We read the code as text at commit \texttt{d07304b} (2023-06-17)
and never ran it. The repository was compromised in October 2026: later commits add code that
collects credentials when the folder is opened in an editor. Readers should not clone it.}, the
test set is the release's test folder, and the validation set is nine training photographs drawn
with NumPy seed~0 from the sorted file names (2, 30, 72, 87, 101, 118, 136, 149, 150), which we
reproduce exactly. Photographs are halved with a Lanczos filter and cut into $128\times128$ patches
on a grid from the origin. The test set is 2,737 such patches, 394 of them with rust; the
validation and test patch counts recomputed from the release match the paper's Table~2 exactly.

Our split is therefore the paper's test set (10 photographs), its validation set without 101
(8 photographs), and its training set without the four twins (77 photographs).

% =================================================================================================
\section{Methods}
\label{sec:methods}

\subsection{Models}
\label{sec:models}

\begin{description}
  \item[UNet (NWRD paper).] The architecture of the NWRD paper: four encoder and four decoder
  blocks of two $3\times3$ convolutions with batch normalisation, 64 to 1024 channels, max-pooling
  and transposed-convolution upsampling, trained from scratch.
  \item[DeepLabV3+ R50.] DeepLabV3+~\cite{chen2018deeplabv3plus} with an ImageNet-pretrained
  ResNet-50 encoder, as implemented in segmentation-models-pytorch 0.5.0~\cite{iakubovskii2019smp}
  with its defaults (output stride 16, ASPP rates 12, 24 and 36, 256 decoder channels).
  \item[CANet-B4.] A context-aggregation network on an ImageNet-pretrained
  EfficientNet-B4~\cite{tan2019efficientnet}. Its 1/32-resolution features (reduced to 512
  channels) feed a context aggregation module (CAM) of a serial dilated branch (dilations 2, 4 and
  8) and three parallel branches that pool by 2, 4 and 8, concatenated and re-weighted by a
  per-pixel, per-channel sigmoid gate. An asymmetric decoder fuses the upsampled context with the
  encoder's 1/4-resolution features.
\end{description}

\note{CANet-B4 was designed in this project; if it draws on a published context-aggregation design,
cite it here.}

\subsection{Windows and scales}
Photographs come in 10 sizes, so a fixed grid either pads (scoring black borders that never occur
in a real photograph as easy background) or drops the edges. Along each side of length $L$ we
instead place $\lceil L/P\rceil$ windows of size $P$ spread evenly from $0$ to $L-P$; neighbouring
windows overlap by at most 46 pixels at 512px, and each overlap is split at its midpoint, so every
pixel is owned by exactly one window and scored exactly once. Nothing is padded and no edge is
dropped. We run two scales: \textbf{512px}, windows of $512\times512$ at native resolution, and
\textbf{128px}, the paper's scale, windows of $128\times128$ on photographs halved with an
anti-aliased Lanczos filter (masks halved by area).

\subsection{Training}
\label{sec:training}
Every model is trained with the same protocol (Table~\ref{tab:training}). The training sample is
every window with rust plus a fixed 20\% of the others, 2,204 windows at 512px and 6,282 at 128px.
The loss is
\begin{equation}
\mathcal{L} = 0.7\,\mathcal{L}_{\mathrm{Dice}} + 0.3\,\mathcal{L}_{\mathrm{BCE}}(w_{+}),
\end{equation}
where the positive weight $w_{+}$ is the ratio of background to rust pixels in the training sample
(63.1 at 512px, 43.9 at 128px).

\begin{table}[t]
\centering
\caption{Training protocol, identical for every model (source: \texttt{docs/protocol.md}, \S4).}
\label{tab:training}
\small
\begin{tabular}{ll}
\toprule
Setting & Value \\
\midrule
Optimiser & AdamW~\cite{loshchilov2019adamw}, learning rate $5\times10^{-5}$, weight decay $10^{-4}$ \\
Schedule & Cosine annealing with warm restarts~\cite{loshchilov2017sgdr} ($T_0=5$, $T_{\mathrm{mult}}=2$), 30 epochs \\
Batch, precision & 8; mixed precision; gradient clipping at norm 1.0 \\
Encoder freezing & Pretrained encoders frozen for the first 3 epochs \\
Augmentation & Flips, $90^\circ$ rotations, brightness and contrast, elastic deformation \\
Model selection & Epoch with the highest pooled validation IoU at threshold 0.5 \\
Threshold & Chosen per run on validation (IoU-maximising, 0.01--0.99), then frozen for test \\
Seeds & 42, 43, 44 \\
Hardware & One NVIDIA Tesla T4 per run (Kaggle) \\
\bottomrule
\end{tabular}
\end{table}

\subsection{Evaluation}
Each model predicts whole test photographs, window by window, and is scored at native resolution
on every pixel once. IoU, Dice (equal to F1), precision and recall of the rust class are computed
from confusion counts pooled over all test pixels,
\begin{equation}
\mathrm{IoU} = \frac{TP}{TP+FP+FN}, \qquad \mathrm{Dice} = \frac{2\,TP}{2\,TP+FP+FN}.
\end{equation}
To compare with the published model, the same predictions are also scored on the paper's 2,737
test patches, in the paper's order and with its per-batch averaging of precision, recall and F1
at threshold 0.5, both on all patches (the paper's Table~3) and on the patches with rust (closest
to its Table~4).

\subsection{Statistical analysis}
\label{sec:stats}
The 10 test photographs are the units of uncertainty: pixels of one photograph share plant, light
and annotator. For each pair of models we report the seed-averaged difference in pooled test IoU
with a 95\% interval from a paired bootstrap~\cite{efron1993bootstrap} that resamples whole test
photographs (5,000 resamples, the same draws for both models and every seed), and a paired $t$-test
over the three seeds, which captures training-run variance but has little power.

The study was pre-registered on 2026-10-22, before any run, with CANet-B4 as the candidate: it
\emph{improves on} a baseline if its mean test IoU is higher, the interval excludes zero and it is
higher on at least two of three seeds; it \emph{is worse} if the interval lies below zero; the two
are \emph{comparable} otherwise. We report those verdicts (Section~\ref{sec:main512}). The
comparisons of DeepLabV3+ against the other models, and the ablation, were added after the 512px
results and are exploratory; they use the same statistics.

\subsection{Ablation}
\label{sec:ablation-design}
To attribute DeepLabV3+'s advantage, we train three variants at 512px, each differing from a
trained model in one factor (Table~\ref{tab:ablation}): DeepLabV3+ R50 without ImageNet weights
(pretraining); a UNet decoder on the same pretrained ResNet-50 encoder (context module: ASPP against
the UNet decoder); and DeepLabV3+ with CANet-B4's EfficientNet-B4 encoder (encoder), which also
compares CAM against ASPP on the same encoder. Training and evaluation are unchanged; encoders
without pretraining are never frozen.

\note{The ablation protocol must be written into \texttt{docs/protocol.md} before its runs. With
three seeds, differences below about 0.03 IoU are within run-to-run variation (see
Section~\ref{sec:discussion}).}

% =================================================================================================
\section{Results}
\label{sec:results}

\subsection{512px}
\label{sec:main512}

DeepLabV3+ R50 has the highest test IoU (Table~\ref{tab:main512}). It is higher than the paper's
UNet by \val{512.dlr-unet.diou} IoU (95\% CI \val{512.dlr-unet.ci}, bootstrap $p$
\val{512.dlr-unet.p}; higher on \val{512.dlr-unet.wins} seeds, seed-level $p$
\val{512.dlr-unet.ttp}) and than CANet-B4 by \val{512.dlr-canet.diou} (95\% CI
\val{512.dlr-canet.ci}). Under the pre-registered rule, CANet-B4 is worse than DeepLabV3+ R50 and
comparable to the UNet ($\Delta$IoU \val{512.canet-unet.diou}, 95\% CI \val{512.canet-unet.ci};
Table~\ref{tab:pairwise512}).

\begin{table}[t]
\centering
\caption{Test results at 512px: the 10 test photographs, every pixel scored once, at each run's
validation-selected threshold (mean $\pm$ sample standard deviation over seeds). ``IoU at 0.5'' is
the IoU at a fixed threshold of 0.5.}
\label{tab:main512}
\small
\resizebox{\linewidth}{!}{\begin{tabular}{lccccccc}
\toprule
Model & Seeds & IoU & Dice & Precision & Recall & IoU at 0.5 & Threshold \\
\midrule
DeepLabV3+ R50 & 3 & \textbf{0.4572\,$\pm$\,0.0084} & 0.6275\,$\pm$\,0.0079 & 0.6078\,$\pm$\,0.0419 & 0.6516\,$\pm$\,0.0319 & 0.4423\,$\pm$\,0.0015 & 0.82 \\
CANet-B4 & 3 & 0.4235\,$\pm$\,0.0072 & 0.5950\,$\pm$\,0.0070 & 0.5779\,$\pm$\,0.0245 & 0.6140\,$\pm$\,0.0120 & 0.4027\,$\pm$\,0.0113 & 0.92 \\
UNet (NWRD paper) & 3 & 0.3983\,$\pm$\,0.0119 & 0.5697\,$\pm$\,0.0121 & 0.5087\,$\pm$\,0.0469 & 0.6543\,$\pm$\,0.0535 & 0.2900\,$\pm$\,0.0176 & 0.96 \\
\bottomrule
\end{tabular}
}
\end{table}

\begin{table}[t]
\centering
\caption{Differences between models at 512px (first minus second). The interval and the first
$p$-value resample test photographs; the second $p$-value is a paired $t$-test over seeds.}
\label{tab:pairwise512}
\small
\resizebox{\linewidth}{!}{\begin{tabular}{llccccc}
\toprule
Comparison & Isolates & $\Delta$IoU & 95\% CI & $p$ (photographs) & $p$ (seeds) & Higher on \\
\midrule
DeepLabV3+ R50 $-$ UNet (NWRD paper) & DeepLabV3+ against the NWRD paper architecture & $+0.0589$ & [$+0.0415$, $+0.0797$] & $<$0.001 & 0.037 & 3 of 3 \\
DeepLabV3+ R50 $-$ CANet-B4 & DeepLabV3+ against the pre-registered candidate & $+0.0337$ & [$+0.0127$, $+0.0476$] & 0.002 & 0.010 & 3 of 3 \\
CANet-B4 $-$ UNet (NWRD paper) & the pre-registered comparison & $+0.0252$ & [$-0.0006$, $+0.0612$] & 0.062 & 0.138 & 3 of 3 \\
\bottomrule
\end{tabular}
}
\end{table}

\subsection{Against the published model, on its own test patches}
On the paper's 2,737 test patches, scored as the paper scores them, DeepLabV3+ R50 reaches F1
\val{512.dlr.pp.f1} against the published \val{paper.t3.f1} (\val{512.dlr.pp.f1.vst3}), and on the
patches with rust F1 \val{512.dlr.pp.f1rust} and IoU \val{512.dlr.pp.iourust} against the published
\val{paper.t4.f1} and \val{paper.t4.iou} (Table~\ref{tab:paper512}). These differences are not
paired tests: the published scores are single runs with no reported variance, come from a model
trained with APF, a focal loss and a longer schedule, and include the leak of
Section~\ref{sec:data}, which favours the published model. The paper's UNet architecture retrained
under our protocol reaches F1 \val{512.unet.pp.f1} on all patches.

\begin{table}[t]
\centering
\caption{The paper's 2,737 test patches (128px patches of the half-size test photographs), scored
with the paper's per-batch averaging at threshold 0.5. Published rows are as reported by Anwar et
al.~\cite{anwar2023nwrd}; our models' predictions at 512px are resized to half size by area.}
\label{tab:paper512}
\small
\begin{tabular}{lcccc}
\toprule
Model & Precision & Recall & F1 & IoU \\
\midrule
\multicolumn{5}{l}{\emph{All test patches, threshold 0.5, the authors' per-batch averaging}} \\
UNet + APF, published (Table 3) & 0.506 & 0.624 & 0.557 & -- \\
DeepLabV3+ R50 & 0.519 & 0.700 & 0.592 & -- \\
CANet-B4 & 0.476 & 0.678 & 0.556 & -- \\
UNet (NWRD paper) & 0.319 & 0.827 & 0.442 & -- \\
\midrule
\multicolumn{5}{l}{\emph{Test patches with rust (closest to Table 4), threshold 0.5}} \\
UNet + APF, published (Table 4) & 0.593 & 0.552 & 0.564 & 0.438 \\
DeepLabV3+ R50 & 0.631 & 0.732 & 0.677 & 0.512 \\
CANet-B4 & 0.592 & 0.716 & 0.648 & 0.479 \\
UNet (NWRD paper) & 0.475 & 0.834 & 0.605 & 0.434 \\
\bottomrule
\end{tabular}

\end{table}

\subsection{Per photograph}
% From generated/tab_photos_512: DeepLabV3+ R50 is highest on 105, 123, 29, 62, 63 and 67.
DeepLabV3+ R50 has the highest IoU on six of the ten test photographs (Table~\ref{tab:photos512}).
\note{Discuss the exceptions: photograph 80, where the UNet is highest (80 is one of the two test
photographs whose twin the paper trained on), and 25, 77 and 89, where CANet-B4 is; and the low
scores on the sparsest photograph, 77 (0.08\% rust).}

\begin{table}[t]
\centering
\caption{Test IoU of each test photograph at 512px, mean over seeds.}
\label{tab:photos512}
\small
\begin{tabular}{llcccc}
\toprule
Photograph & Camera & Rust & DeepLabV3+ R50 & CANet-B4 & UNet (NWRD paper) \\
\midrule
105 & PAR-LX1M & 0.57\% & 0.3857 & 0.3606 & 0.2603 \\
123 & PAR-LX1M & 0.97\% & 0.4708 & 0.4513 & 0.4329 \\
25 & NIKON D3200 & 0.32\% & 0.3980 & 0.4100 & 0.2942 \\
29 & NIKON D3200 & 0.35\% & 0.5257 & 0.4948 & 0.4661 \\
62 & NIKON D3200 & 1.04\% & 0.4363 & 0.3959 & 0.3760 \\
63 & NIKON D3200 & 1.24\% & 0.4875 & 0.4369 & 0.4493 \\
67 & PAR-LX1M & 0.76\% & 0.6136 & 0.5373 & 0.5772 \\
77 & PAR-LX1M & 0.08\% & 0.2226 & 0.2358 & 0.1719 \\
80 & PAR-LX1M & 0.25\% & 0.3051 & 0.2534 & 0.3613 \\
89 & PAR-LX1M & 0.47\% & 0.3482 & 0.3592 & 0.2995 \\
\bottomrule
\end{tabular}

\end{table}

\begin{figure}[t]
\centering
\figureorpending{training_curves_512px.png}{0.85\linewidth}
\caption{Training at 512px. Left: validation loss on windows; right: validation IoU on whole
photographs at threshold 0.5. Lines are means over seeds, bands span the seeds' minimum and
maximum.}
\label{fig:curves512}
\end{figure}

\subsection{128px, the paper's scale}
\label{sec:main128}
\pending{128px results: Table~\ref{tab:main128}, the pairwise differences and the paper-patch
comparison at the paper's own scale.}

\begin{table}[t]
\centering
\caption{Test results at 128px (photographs halved, $128\times128$ windows).}
\label{tab:main128}
\small
\resizebox{\linewidth}{!}{\begin{tabular}{lccccccc}
\toprule
Model & Seeds & IoU & Dice & Precision & Recall & IoU at 0.5 & Threshold \\
\midrule
DeepLabV3+ R50 & 3 & \textbf{0.4098\,$\pm$\,0.0150} & 0.5812\,$\pm$\,0.0151 & 0.5255\,$\pm$\,0.0196 & 0.6504\,$\pm$\,0.0101 & 0.3784\,$\pm$\,0.0098 & 0.89 \\
CANet-B4 & 3 & 0.4073\,$\pm$\,0.0127 & 0.5788\,$\pm$\,0.0129 & 0.5448\,$\pm$\,0.0137 & 0.6173\,$\pm$\,0.0119 & 0.3673\,$\pm$\,0.0111 & 0.91 \\
UNet (NWRD paper) & 3 & 0.4011\,$\pm$\,0.0365 & 0.5720\,$\pm$\,0.0367 & 0.4914\,$\pm$\,0.0595 & 0.6888\,$\pm$\,0.0103 & 0.3426\,$\pm$\,0.0403 & 0.88 \\
\bottomrule
\end{tabular}
}
\end{table}

\subsection{Ablation}
\label{sec:ablation}
\pending{the ablation runs. Report which single change accounts for most of DeepLabV3+'s advantage
over the UNet, with each difference's interval (Table~\ref{tab:ablation}).}

\begin{table}[t]
\centering
\caption{Ablation at 512px. Each variant differs from a trained model in one factor; differences
are against DeepLabV3+ R50 with 95\% intervals over test photographs.}
\label{tab:ablation}
\small
\resizebox{\linewidth}{!}{\begin{tabular}{llllcc}
\toprule
Model & Encoder & Pretraining & Context module & Test IoU & $\Delta$IoU vs DeepLabV3+ R50 [95\% CI] \\
\midrule
DeepLabV3+ R50 & ResNet-50 & ImageNet & ASPP & 0.4572\,$\pm$\,0.0084 & reference \\
CANet-B4 & EfficientNet-B4 & ImageNet & CAM & 0.4235\,$\pm$\,0.0072 & $-0.0337$ [$-0.0476$, $-0.0127$] \\
UNet (NWRD paper) & plain convolutions & none & UNet decoder & 0.3983\,$\pm$\,0.0119 & $-0.0589$ [$-0.0797$, $-0.0415$] \\
DeepLabV3+ R50, no pretraining & ResNet-50 & none & ASPP & \textcolor{red}{pending} & \textcolor{red}{pending} \\
UNet R50 & ResNet-50 & ImageNet & UNet decoder & \textcolor{red}{pending} & \textcolor{red}{pending} \\
DeepLabV3+ B4 & EfficientNet-B4 & ImageNet & ASPP & \textcolor{red}{pending} & \textcolor{red}{pending} \\
\bottomrule
\end{tabular}
}
\end{table}

% =================================================================================================
\section{Discussion}
\label{sec:discussion}

\paragraph{Why DeepLabV3+ might win.} \pending{to be written from the ablation.} Plausible
explanations, each tested by one variant, are ImageNet pretraining (77 training photographs are
few for a from-scratch encoder), the wider receptive field of ASPP on sparse, elongated rust
stripes, and the ResNet-50 encoder itself.

\paragraph{Scores against annotation agreement.} The best test IoU at 512px,
\val{512.dlr.iou}, is close to the median agreement between two annotations of the same photograph
(\val{data.twins.median}). Absolute IoUs on NWRD should therefore be read against this agreement
rather than against scores on datasets with denser or more consistent annotation.

\paragraph{Run-to-run variation.} \note{A second, accidental run of DeepLabV3+ R50 at 512px with
identical code and seeds (Kaggle version 3, in \texttt{runs\_kaggle/}, not in \texttt{results/})
gave per-seed test IoUs differing from the reported run by up to about 0.03, most likely from
non-deterministic GPU convolution algorithms (\texttt{cudnn.benchmark}). The DeepLabV3+-against-CANet
verdict holds against it. Collect it into \texttt{results/} as a replicate before citing its
numbers, and disclose it here.}

\paragraph{Limitations.} APF and the paper's focal loss and schedule are not reproduced: we test the
paper's architecture under our protocol, and its full method only through its published numbers.
The test set is 10 photographs from two cameras, so intervals are wide. Three seeds give little
power to the seed-level test. The ablation is exploratory and varies one factor at a time, so it
does not measure interactions.

% =================================================================================================
\section{Conclusion}
\label{sec:conclusion}
On the public NWRD release, under a protocol free of the release's cross-split duplicates, an
ImageNet-pretrained DeepLabV3+ with a ResNet-50 encoder segments stripe rust better than the NWRD
paper's UNet retrained under the same protocol (test IoU \val{512.dlr.iou.pm} against
\val{512.unet.iou.pm} at 512px), and better than the published UNet with APF on the paper's own test
patches (F1 \val{512.dlr.pp.f1} against \val{paper.t3.f1}). \pending{128px and the ablation's
answer.} The disagreement between independent annotations of the same photographs, at a median IoU
of \val{data.twins.median}, suggests that further gains on this dataset need better annotation as
much as better models.

\section*{Data and code availability}
The data are the public NWRD release (Zenodo, doi:10.5281/zenodo.21672928). \note{State where the
code and the result files will be released.}

\bibliographystyle{unsrtnat}
\bibliography{references}

\appendix
\section{Where each number comes from}
\label{app:trace}
Every model number in this paper is generated by \texttt{scripts/build\_paper.py} from the study's
result files: per-model summaries from \texttt{results/<scale>/comparison.json}; differences
between models from \texttt{results/<scale>/pairwise.json} (the same photograph bootstrap as the
pre-registered rule); per-photograph scores from each run's \texttt{test\_counts.npz}; the
published scores from the NWRD paper's Tables~3 and~4; annotation agreement from
\texttt{results/dataset/duplicates.json}. Dataset facts written in the text come from
\texttt{docs/dataset.md} and \texttt{results/dataset/summary.md}.

\end{document}
